% ============================================================================
% Fichier cible : ARIEL_LOGMINER_MEMOIRE_FINAL/main.tex, chapitre étoilé « Résumé »
% Lignes à remplacer : 278 à 382 (de « Ce mémoire présente Ariel \logminer{}, une
%   architecture multi-agents légère » jusqu'à « pour l'analyse de journaux
%   hétérogènes. »). Conserver les lignes 275-276 (\chapter*, \addcontentsline)
%   et le bloc des mots-clés (lignes 384-390).
% Auteur : Agent 1 (intégrité factuelle), chantier 2026-11-05 ; paragraphe
%   « hypothèses » réécrit par l'Agent 2 (même date). VERSION À INTÉGRER À LA
%   PLACE de agent_1/sections/main.tex__resume_fr.tex (décision C7).
% Modification de l'Agent 2 : seul l'avant-dernier paragraphe (« Aucune des
%   cinq hypothèses n'est validée [...] ») change ; il reprend les statuts du
%   tableau 7.1 de agent_2/sections/chapitre6_conclusion.tex__statut_hypotheses_tableau_7.1.tex.
%   La phrase « la question méthodologique (QR6) est la mieux couverte » est
%   retirée (L2 §2.1) ; une phrase sur l'absence d'analyse conjointe est
%   ajoutée (L3, non-recevabilité). Aucun chiffre nouveau (règle C8) ; toutes
%   les affirmations figurent dans le corps : §4.2.3 (tout positif), §4.4
%   (règle de sévérité), §4.6.2 (règle de port), §4.3 (HDFS), §4.9 (mémoire),
%   §4.9.1 (idempotence transposable), §4.8 (schéma commun), tableau 7.1.
% Statut : définitif pour le paragraphe « hypothèses » ; les paragraphes sur
%   l'architecture (benchmark, multi-VM) restent ceux de l'Agent 1, corrigés sur
%   l'exactitude factuelle seulement ; l'Agent 5 part de CE fichier pour les
%   reformuler au fond.
% Sources : voir LIVRABLE_1_integrite.md, §4, et LIVRABLE_2_hypotheses.md, §4.
% ============================================================================

Ce mémoire présente Ariel \logminer{}, une architecture multi-agents légère
destinée à l'analyse de journaux système et réseau hétérogènes, dont la
distribution est évaluée dans une topologie de laboratoire. Un agent est ici
une entité logicielle dotée de capacités, d'un état, d'une mémoire et d'une
politique locale. Lorsqu'un objectif borné lui est confié, il peut évaluer les
tâches disponibles, proposer leur exécution ou les refuser, exécuter des
handlers compatibles, enchaîner des sous-tâches, reprendre un état persistant
et déléguer une opération exigeant une autre capacité. La négociation repose
sur un protocole Contract Net explicite et sur un bus local ou Redis Streams.
L'allocation et la distribution sont évaluées par des campagnes en laboratoire,
tandis que la poursuite bornée d'objectifs est validée fonctionnellement.
L'intelligence générale, la définition autonome des finalités et le
fonctionnement d'un SOC industriel restent hors périmètre.

L'évaluation met particulièrement en évidence l'influence du protocole
expérimental sur les performances. Sur CICIDS2017, un Random Forest évalué
avec une séparation aléatoire stratifiée obtient un F1 moyen de
$0{,}995142$ sur cinq graines. Lorsque chacun des cinq scénarios d'attaque
est tenu hors entraînement, le F1 moyen tombe à $0{,}157744$. Parmi cinq
modèles légers comparés sous ce holdout, la régression logistique obtient
le meilleur F1 moyen, $0{,}233670$, tandis que HistGradientBoosting obtient
une PR-AUC macro de $0{,}569002$ pour un F1 presque nul au seuil retenu. Un
changement de modèle modifie donc certaines métriques sans supprimer la
fragilité inter-scénarios.

Pour HDFS et BGL, le protocole sépare strictement entraînement, validation et
test. Drain3 est ajusté uniquement sur l'entraînement, puis gelé avant les
deux partitions suivantes. Au niveau \texttt{block\_id}, Histogram atteint un
F1 de $0{,}868373$ (IC à 95~\% : $[0{,}856771\,;\,0{,}879493]$) sur les
$115\,013$ blocs de test HDFS, dont $1\,679$ positifs, avec un rappel de $1$
et une précision de $0{,}767367$, identiques pour cinq graines. Sur BGL, le F1
vaut $0{,}913698$, mais $89{,}8081\,\%$ des événements de test portent un
template inconnu de l'entraînement, et la décision d'Histogram coïncide
exactement avec la règle « sévérité absente de l'entraînement normal ». Ce
score mesure donc une nouveauté de sévérité plutôt qu'une modélisation de la
dynamique des journaux ; la règle « template inconnu » n'atteint que
$0{,}277885$.

Sur CSE-CIC-IDS2018, une séparation entre le 15 et le 16 février donne un F1
moyen de $0{,}999212$ à la régression logistique contre $0{,}401375$ au
Random Forest. L'audit montre cependant que \texttt{Dst Port} seul atteint
$0{,}999720$, score qu'une règle fixe sur le port de destination (inférieur à
1024) obtient également, et que le transfert inverse, du 16 vers le
15 février, ne dépasse pas $0{,}234195$. Le score principal repose donc sur
un raccourci propre à la journée de test.

La validation multiformat lit, parse et normalise $7\,001/7\,001$ unités, dont
une seule fixture Apache. Ce résultat mesure la couverture fonctionnelle des
formats testés ; le remplissage du schéma commun varie de sept champs standard
sur sept (syslog, BGL) à aucun (flux réseau tabulaires). Une politique
open-set calibrée uniquement sur les sources connues, puis gelée, rejette
$20/20$ fichiers inconnus et écarte à tort $3/46$ fichiers connus ; les
43 fichiers connus acceptés sont correctement orientés et la marge de décision
n'est pas une probabilité calibrée. Aucune comparaison avec une chaîne sans
routage n'a été exécutée : le bénéfice prédictif du routage spécialisé n'est
pas établi.

Le benchmark contrôlé couvre quatre architectures, quatre charges et dix
répétitions, soit $160$ exécutions et $264\,000$ traitements sans échec. Sur
les microtâches du protocole, le monolithe domine et les architectures à
agents ajoutent un coût d'orchestration. Lors d'une panne appariée survenant
après persistance, les architectures à agents ne dupliquent aucun effet,
contrairement au monolithe et aux workers centralisés ; doter ces derniers du
même magasin d'idempotence supprime également les doublons.

Une campagne d'endurance d'une heure, sur deux machines virtuelles Debian et
Ubuntu reliées à un Redis centralisé sur l'hôte du laboratoire, mène à terme
$1\,800$ tâches sur $1\,800$, réparties 451 et $1\,349$ selon les profils de
capacité configurés ; un arrêt contrôlé après persistance conduit à une
réattribution et à un rejeu sans effet dupliqué. En saturation, la capacité
de cette topologie se situe entre 8 et 16 tâches/s, et la seconde machine ne
traite que $1{,}23\,\%$ des tâches réussies : elle n'augmente pas la capacité
observée. La haute disponibilité n'a pas été évaluée.

Le corrélateur, évalué sur un corpus synthétique contrôlé, obtient une
précision de $0{,}764706$, un rappel de $0{,}866667$ et un F1 par paires de
$0{,}812500$. Sur des corpus générés plus grands, son rappel tombe à
$0{,}5973$ ; lorsque 30~\% des incidents enchaînent plusieurs catégories
d'événements, une variante réduite à l'hôte et à la fenêtre temporelle obtient
un meilleur F1 ($0{,}7268$ contre $0{,}6496$).

Dans la campagne CNP multi-source, la condition M exécute réellement les
artefacts compatibles sur $1400/1401$ unités, avec un seul fallback Apache et
aucune erreur. Sa latence moyenne est de $0{,}105420$~s, contre $0{,}031027$~s
pour la condition heuristique H. L'absence de labels homogènes exclut le
calcul d'une exactitude globale multi-source ; le F1 Linux/auth de
$0{,}390244$ reste descriptif de la source sélectionnée.

Le dépôt contient aussi une procédure de comparaison des modèles et un mode
dry-run. Les tests couvrent certaines branches, mais pas un cycle complet et
répété de promotion et de rejet. Le bénéfice global de cette procédure n'a pas
été évalué.

Aucune des cinq hypothèses n'est validée. H1 et H2 sont partiellement
soutenues, mais ce qu'elles établissent relève d'exigences de réalisation :
chaque unité est lue et émise dans le schéma commun, et le magasin
d'idempotence supprime les effets dupliqués quelle que soit l'architecture.
H3 reste non tranchée, faute de comparaison avec une chaîne sans routage. H4
n'est pas tranchée : sous séparation temporelle pour CICIDS2017, sur
CSE-CIC-IDS2018 et sur BGL, aucun modèle ne dépasse la règle triviale évaluée
sur le même test, aucune baseline n'est évaluée au niveau bloc sur HDFS, et
aucun budget de ressources n'a été fixé. H5 est réfutée selon son propre
critère : la mémoire n'apporte aucun gain établi, et la priorisation qu'elle
devait améliorer n'a pas été mesurée. Aucune expérience n'analyse
conjointement un même incident à partir de journaux système et de flux
réseau. La contribution principale d'Ariel \logminer{} réside ainsi dans
l'intégration architecturale et méthodologique d'une chaîne locale, modulaire
et expérimentalement traçable pour l'analyse de journaux hétérogènes.
